% ============================================================
% Chapter 3 — Data
% ============================================================

\chapter{Data}
\label{chap:data}

% ────────────────────────────────────────────────────────────
\section{Data Collection}
\label{sec:data-collection}

The dataset originates from Danske Bank's retail payment system under a data sharing agreement. A \emph{transaction table} contains \todo{FILL} retail payments recorded over \todo{FILL: date range}, each carrying operational metadata: sending and receiving accounts, amount, currency, payment method, channel, timestamp, and institutional routing codes. A separate \emph{fraud labels table} lists confirmed fraud cases identified through the bank's investigation process. The two tables were joined on a matched transaction identifier, yielding a binary label column \texttt{CONFIRMED\_RISK}.

\paragraph{Label semantics.}
A label of 0 denotes \emph{absence from the known fraud table}, not confirmed legitimacy. The dataset therefore contains an unknown proportion of false negatives. This motivates the use of PR-AUC as the primary metric rather than accuracy or AUROC, both of which treat label 0 as a reliable negative. The implications for evaluation are discussed further in \Cref{sec:eval-framework}.

% ────────────────────────────────────────────────────────────
\section{Dataset Overview}
\label{sec:data-overview}

\Cref{tab:dataset-overview} summarises the working dataset.

\begin{table}[h]
    \centering
    \caption{Dataset overview.}
    \label{tab:dataset-overview}
    \begin{tabular}{ll}
        \toprule
        \textbf{Property} & \textbf{Value} \\
        \midrule
        Total transactions              & \tabletodo{FILL} \\
        Time period                     & \tabletodo{FILL} \\
        Confirmed fraud cases           & \tabletodo{FILL (\% rate)} \\
        Unique sending accounts         & \tabletodo{FILL} \\
        Unique receiving accounts       & \tabletodo{FILL} \\
        \bottomrule
    \end{tabular}
\end{table}

% ────────────────────────────────────────────────────────────
\section{Feature Engineering}
\label{sec:feature-engineering}

\subsection{Transaction-Level Features}

Each transaction carries the following engineered features:
\begin{itemize}[nosep]
  \item $\log(1 + \texttt{BASEVALUE})$, z-score normalised
  \item One-hot encoding of \texttt{CHANNEL}
  \item One-hot encoding of \texttt{PAYMENTSUBMETHOD}
  \item Binary clearing express flag
  \item Binary international flag
  \item One-hot encoding of \texttt{ACCOUNTBRANCH\_TBE}
  \item Cyclical time encoding: $\sin/\cos$ of hour-of-day and day-of-week (4 dimensions)
\end{itemize}

All normalisation parameters and one-hot vocabularies are fitted exclusively on the training split.

\subsection{Node Features}

Node features are computed by aggregating training-period transactions only, preventing information leakage from future data.

\paragraph{InternalAccount.}
\todo{FILL: list current features (out-degree, amount stats, counterparty diversity, channel diversity, time behaviour)}

\paragraph{ExternalAccount.}
\todo{FILL: list current features (in-degree, received amount stats, sender diversity, sender bank diversity)}

% ────────────────────────────────────────────────────────────
\section{Temporal Split}
\label{sec:temporal-split}

Transactions are split chronologically to mirror deployment conditions:
\begin{itemize}[nosep]
  \item \textbf{Training}: before \todo{FILL: first cutoff}
  \item \textbf{Validation}: \todo{FILL: first cutoff} to \todo{FILL: second cutoff}
  \item \textbf{Test}: after \todo{FILL: second cutoff}
\end{itemize}

This ensures the model is always evaluated on transactions that postdate its training window, testing generalisation to potentially unseen accounts and counterparties.

% ────────────────────────────────────────────────────────────
\section{Class Imbalance}
\label{sec:class-imbalance}

The fraud rate of \todo{FILL}\% constitutes severe class imbalance. This is addressed via inverse-frequency weighting of the binary cross-entropy loss:
\[
  w_{\text{pos}} = \frac{n_{\text{neg}}}{n_{\text{pos}}}
\]
The weight is recomputed for each training split.
